\documentclass[12pt,a4paper]{article}

\usepackage{iftex}
\ifPDFTeX
  \usepackage[T2A]{fontenc}
  \usepackage[utf8]{inputenc}
\else
  \usepackage{fontspec}
  \IfFontExistsTF{Times New Roman}
    {\setmainfont{Times New Roman}}
    {\setmainfont{DejaVu Serif}}
  \IfFontExistsTF{Consolas}
    {\setmonofont{Consolas}}
    {\setmonofont{DejaVu Sans Mono}}
\fi
\usepackage[ukrainian]{babel}
\usepackage{microtype}
\usepackage{geometry}
\usepackage{amsmath,amssymb}
\usepackage{array}
\usepackage{booktabs}
\usepackage{enumitem}
\usepackage{longtable}
\usepackage{titlesec}

\geometry{left=25mm,right=20mm,top=20mm,bottom=22mm}
\setlength{\parindent}{1.25cm}
\setlength{\parskip}{0.25em}
\setlist{itemsep=0.2em,topsep=0.25em}
\pagestyle{plain}
\allowdisplaybreaks

\titleformat{\section}{\large\bfseries}{\thesection.}{0.55em}{}
\titleformat{\subsection}{\normalsize\bfseries}{\thesubsection.}{0.5em}{}

\newcommand{\code}[1]{\texttt{#1}}

\begin{document}

\begin{titlepage}
\begin{center}
  Міністерство освіти і науки України\\
  Київський національний університет імені Тараса Шевченка\\[5em]

  {\Large\bfseries КУРСОВИЙ ПРОЄКТ}\\[1em]
  {\large з дисципліни «Теорія програмування»}\\[2em]

  {\LARGE\bfseries SIPL Language Toolkit}\\[0.7em]
  {\large Парсер, інтерпретатор, візуалізація та семантичні терми}\\[5em]
\end{center}

\begin{flushright}
  Виконав:\\
  студент групи ТТП-32\\
  Кашинський Владислав
\end{flushright}

\vfill
\begin{center}
  Київ - 2026
\end{center}
\end{titlepage}

\tableofcontents
\newpage

\section{Мета та постановка задачі}

Мета курсового проєкту - створити розширюваний набір інструментів для
навчальної мови SIPL, у якому граматика, парсер, AST, виконання,
візуалізація та семантика працюють через стабільні програмні інтерфейси.
Реалізація не прив'язана до одного варіанта практичної роботи й підтримує
цілий клас SIPL-програм.

Для виконання обрано модулі, наведені в таблиці~\ref{tab:tasks}.

\begin{table}[h]
\centering
\caption{Заявлені модулі проєкту}
\label{tab:tasks}
\begin{tabular}{>{\raggedright\arraybackslash}p{18mm}p{100mm}r}
\toprule
ID & Зміст & Бали\\
\midrule
CORE-01 & Машинозчитувана граматика, позитивні та негативні тести & 0.5\\
CORE-02 & Парсер і типізоване AST із позиціями у тексті & 1.5\\
CORE-03 & Інтерпретатор SIPL із лімітом кроків & 2.0\\
CORE-04 & Подієва модель виконання та trace & 1.0\\
CORE-05 & CLI та структурований JSON-режим & 1.0\\
VIS-01 & Parse tree у форматах Text, JSON і DOT & 1.5\\
VIS-02 & AST та порівняння Parse Tree $\leftrightarrow$ AST & 1.0\\
SEM-01 & Автоматична побудова семантичного терму & 2.0\\
\midrule
 & Заявлена сума & 10.5\\
\bottomrule
\end{tabular}
\end{table}

За правилами курсового проєкту підсумкова оцінка обмежується 10 балами.

\section{Архітектура}

Інструментарій реалізовано мовою Python. Опис мови зберігається у формальній
граматиці Lark. На її основі парсер будує конкретне синтаксичне дерево, яке
потім перетворюється у власне типізоване AST. Саме AST передається
інтерпретатору, модулю семантичних термів і засобам візуалізації.

Для роботи з усіма можливостями передбачено єдиний командний інтерфейс.
Окремо реалізовано локальну вебсторінку, що використовує ті самі модулі й не
дублює логіку синтаксичного аналізу або виконання програми.

Потік даних має вигляд
\[
\text{Source}
\longrightarrow \text{Parse Tree}
\longrightarrow \text{AST}
\longrightarrow
\begin{cases}
\text{Interpreter}\longrightarrow\text{State + Trace},\\
\text{Semantic Builder}\longrightarrow\text{Semantic Term},\\
\text{Visualizer}\longrightarrow\text{Text/JSON/DOT}.
\end{cases}
\]

Граматика визначена в одному файлі. Інтерпретатор, семантичний модуль та
візуалізатори не аналізують початковий текст повторно, а використовують
спільні AST-вузли. Це усуває дублювання логіки й дозволяє надалі під'єднати
IDE або debugger.

\section{Формальна граматика SIPL}

Програма обмежується ключовими словами \code{begin} та \code{end}.
Послідовність операторів розділяється крапкою з комою. Підтримано
присвоювання, блок, умовний оператор, цикл і \code{skip}.

\begin{verbatim}
program: "begin" stmt_list? "end"
stmt_list: statement (";" statement)* ";"?

statement: assignment | if_stmt | while_stmt | block | "skip"
assignment: NAME ":=" expr
if_stmt: "if" bool_expr "then" statement "else" statement
while_stmt: "while" bool_expr "do" statement
\end{verbatim}

Для арифметики реалізовано операції \code{+}, \code{-}, \code{*},
\code{div}, \code{mod}; для умов - порівняння, \code{not},
\code{and}, \code{or}. Граматика задає пріоритет і ліву асоціативність
операторів. Однорядкові та блокові коментарі ігноруються лексером.

Набір тестів містить 11 коректних і 11 некоректних синтаксичних прикладів,
що перевищує мінімальний критерій CORE-01.

\section{Типізоване AST та SourceSpan}

Парсер Lark спочатку будує конкретне дерево, після чого модуль
\code{parser.py} створює власні незмінні AST-вузли. Класи \code{Number} і
\code{Variable} представляють числа та змінні. Операціям відповідають
\code{UnaryExpr}, \code{BinaryExpr}, \code{Comparison} і \code{BoolBinary}.
Оператори подано класами \code{Assignment}, \code{Skip}, \code{If},
\code{While} та \code{Block}. Усе дерево об'єднує вузол \code{Program}.

Кожен вузол має \code{SourceSpan}: початковий і кінцевий рядок, колонку та
зміщення у вихідному тексті. Завдяки цьому помилки й події можна точно
пов'язати з фрагментом програми. AST серіалізується у стабільний JSON,
незалежний від внутрішнього формату бібліотеки Lark.

\section{Інтерпретатор і подієва модель}

Стан SIPL-програми подано відображенням
\[
  State = Var\to Int,
\]
реалізованим словником Python. Присвоювання оновлює компоненту стану,
послідовність виконує оператори зліва направо, \code{if} обирає одну гілку,
а \code{while} повторює тіло, доки умова істинна.

Читання неініціалізованої змінної й ділення на нуль породжують зрозумілу
runtime-помилку з позицією у програмі. Ліміт кроків перериває потенційно
нескінченний цикл.

Перед виконанням оператора створюється \code{BeforeStatement}. Після нього
додається \code{AfterStatement}. Зміну значення фіксує
\code{VariableChanged}, а перевірку умови в \code{if} або \code{while}
позначає \code{ConditionEvaluated}. Виконання завершує подія
\code{ProgramFinished}, що містить кінцевий стан.

Кожна подія має \code{SourceSpan}. Trace доступний як об'єкт Python,
текст CLI та JSON, тому його можна повторно використати у debugger або
візуалізаторі.

\section{Семантичні терми}

Семантичний терм будується рекурсивно з AST. Наприклад, для програми
перевірки парності отримано:

\begin{multline*}
WH(S^2(gr,N^{\Rightarrow},1),
AS^N(S^2(sub,N^{\Rightarrow},2)))\;\cdot\\
IF(S^2(eq,N^{\Rightarrow},1),AS^R(1),AS^R(0)).
\end{multline*}

Тут $S^2$ - двомісна суперпозиція, $N^{\Rightarrow}$ - розіменування,
$AS^N$ - присвоювання, $WH$ - цикл, $IF$ - розгалуження, а $\cdot$ -
послідовна композиція. Отже, позначення відповідають матеріалам курсу й
попереднім практичним роботам.

\section{Візуалізація}

Parse tree та AST можна вивести як звичайне текстове дерево з відступами й
позиціями вузлів. Для програмного опрацювання передбачено JSON, а для
графічного подання - формат Graphviz DOT, з якого будуються SVG або PNG.

Окрема команда порівнює кількість вузлів. Parse tree зберігає правила
граматики й токени, тоді як AST вилучає пунктуацію та проміжні синтаксичні
категорії, залишаючи лише семантично значущі конструкції.

\section{Командний інтерфейс}

Реалізовано команди:

\begin{longtable}{p{28mm}p{105mm}}
\toprule
Команда & Призначення\\
\midrule
\endhead
\code{check} & перевірка синтаксичної коректності\\
\code{parse} & побудова типізованого AST\\
\code{run} & виконання та виведення кінцевого стану\\
\code{trace} & покроковий журнал виконання\\
\code{semantic} & побудова семантичного терму\\
\code{tree} & візуалізація parse tree\\
\code{ast} & візуалізація AST\\
\code{compare} & порівняння конкретного й абстрактного синтаксису\\
\bottomrule
\end{longtable}

JSON-відповіді мають поле \code{schema\_version = "1.0"}. Код завершення
0 означає успіх, 2 - синтаксичну або I/O-помилку, 3 - runtime-помилку.

Локальний вебінтерфейс запускається командою \code{sipl-web}.
\par\noindent Адреса сайту:
\code{127.0.0.1:8000}. Він дає змогу ввести програму й стан, виконати
аналіз, переглянути AST, семантичний терм і trace. Зовнішні сервіси не потрібні.

\section{Тестування та демонстрація}

Повний автоматичний набір містить 44 тести. Результат:

\begin{verbatim}
44 passed
TOTAL: 670 statements, 97 missed, coverage 84%
\end{verbatim}

Покриття перевищує вимогу 80\% для ядра. Перевірено синтаксис,
пріоритети операцій, SourceSpan, виконання, коротке замикання умов,
помилки, ліміт кроків, trace, CLI, JSON, DOT і семантичні терми.

Для демонстрації підготовлено три програми. Піднесення до степеня при
$X=2$ і $Y=5$ дає $R=32$. Алгоритм Евкліда для $M=48$ і $N=18$
завершується зі значеннями $M=N=6$. У програмі перевірки парності початкове
$N=5$ переходить у $N=1$, після чого результатом стає $R=1$.

Приклади команд:

\begin{verbatim}
sipl run examples/power.sipl --state X=2 --state Y=5
sipl run examples/gcd.sipl --state M=48 --state N=18
sipl trace examples/parity.sipl --state N=5 --json
sipl semantic examples/parity.sipl
pytest --cov=sipl --cov-report=term-missing
\end{verbatim}

\section{Відтворюваність та обмеження}

Залежності й метадані зафіксовано в \code{pyproject.toml}. Проєкт можна
встановити командою \code{pip install -e ".[dev]"}, після чого доступна
команда \code{sipl}. README містить команди запуску, тестування та опис
кодів завершення.

Поточна версія не реалізує статичну типізацію, IDE, debugger і символічне
доведення коректності. DOT генерується без вбудованого Graphviz-рендерера.
Ці обмеження явно задокументовано й вони не впливають на заявлені модулі.

\enlargethispage{1\baselineskip}
\section{Висновок}

Створено працездатний SIPL Language Toolkit: граматика є єдиним джерелом
синтаксису, парсер будує типізоване AST з позиціями, інтерпретатор повертає
стан і trace, а візуалізація та семантика використовують AST API.

Реалізацію перевірено 44 тестами з покриттям 84\%. Підготовлено приклади
для захисту й відтворюване встановлення. Виконано завдання на максимальні
10 балів за умовами проєкту.

\vspace{0.2em}
\noindent\textbf{Використані матеріали}
\addcontentsline{toc}{section}{Використані матеріали}

\begingroup
\footnotesize
\noindent
1. М.~С.~Нікітченко. \emph{Теорія програмування}. 2020.\quad
2. О.~Шишацька, А.~Криволап, А.~Шишацький. \emph{Теорія програмування. Практикум}. 2024.\quad
3. Технічне завдання \emph{SIPL Language Toolkit}, 2026.\quad
4. Документація бібліотеки Lark Parser.
\endgroup

\end{document}

